\documentclass[10pt,a4paper]{amsart}
\usepackage[margin=19mm]{geometry}
\usepackage{amsmath,amssymb,mathtools}
\usepackage[scr=rsfs,cal=euler]{mathalfa}
\usepackage{xfrac}
\usepackage[unicode,hidelinks,bookmarks=false]{hyperref}
\newcommand{\ie}{i.e.\ }
\newcommand{\cf}{cf.\ }
\newcommand{\resp}{respectively\ }
\newcommand{\Subsection}{\subsection}
\providecommand{\nobreakdash}{\nobreak\hbox{-}}
\setlength{\emergencystretch}{2em}
\multlinegap0pt
\usepackage{mathtools}
\usepackage{amssymb}
\usepackage{tikz}
\newtheorem{lemma}{Lemma}
\title{Polynomiality of the Generalized Verschiebung Degree}
\author{Siqing Zhang}
\date{Source: arXiv:2606.26070v1 (CC BY 4.0)}
\begin{document}
\maketitle

\noindent\emph{Source and licence.} Polynomiality of the Generalized Verschiebung Degree. Version arXiv:2606.26070v1 (CC BY 4.0), distributed by arXiv under the Creative Commons Attribution 4.0 International licence, http://creativecommons.org/licenses/by/4.0/, as recorded in the arXiv licence metadata for this exact version and shown on the abstract page https://arxiv.org/abs/2606.26070v1. Full licence notice: \texttt{licenses/arxiv-2606.26070-CC-BY-4.0.txt}.\par\smallskip

\noindent\textit{Editorial scope.} Counted unit: the article's lemma fac (source0.tex lines 207--212). The article's complete original proof of it is delivered with it (source0.tex lines 213--332, one \texttt{proof} environment of 120 lines). No mathematics is written, repaired or re-derived: the statement and the proof are the article's own lines, in the article's own notation, and no retained line is substituted or re-flowed.  Retained uncounted as the source-local context the proof needs: lines 156--158; lines 160--162; lines 164--166; lines 173--179; lines 198--203.  Nothing was omitted from any retained span, so the package carries no omission record.  No retained line contains a citation, so the wrapper carries no bibliography and no bibliography entry.  No retained line contains a figure environment, an image-inclusion command or drawing code, so no image is delivered and the package declares no asset dependency.  Editorial formatting only, in addition: an AMS \texttt{amsart} preamble that restates the article's own declarations and macros used by the retained lines, namely the environments \texttt{lemma} on separate counters, together with the standard packages those lines need.  The retained environments print in this document as Lemma 1 in order of appearance, whereas the article prints the same items with the numbering of its own full document.  The complete native source of the article as distributed in the arXiv e-print for this exact version is bundled unchanged as \texttt{sources/203665/source0.tex}.
\begin{equation}\label{eq: upper}
    a_1+a_2+a_3\leq P^N-2,\quad |a_2-a_3|\leq a_1\leq a_2+a_3,
\end{equation}

\begin{equation}\label{eq: am}
    a_{i,M}\in \{[a_i]_M,\, P^M-1-[a_i]_M\}
\end{equation}

\begin{equation}\label{eq: lower}
a_{1,M}+a_{2,M}+a_{3,M}\leq P^M-2,\quad |a_{2,M}-a_{3,M}|\leq a_{1,M}\leq a_{2,M}+a_{3,M}.
\end{equation}

\begin{lemma}\label{eqn: four ineq}
    The condition \eqref{eq: upper} is equivalent to the following four inequalities
\begin{align*}
    2P^N -(2a_1+1)-(2a_2+1)-(2a_3+1) &>0,\\
    -(2a_i+1)+\sum_{j\neq i}(2a_j+1)&>0, \quad i=1,2,3.
\end{align*}
\end{lemma}

\begin{equation}\label{eq: Phinb}
    \Phi(n,b)=\begin{cases}
        \frac{n}{2}P^{N-1}+b, & n\ \text{even},\\
       \frac{n+1}{2} P^{N-1}-1-b, & n\ \text{odd}.
        \end{cases}
\end{equation}

\begin{lemma}[The local factorization lemma]\label{lm: local fac}
    Assume that $P\geq 3$ is odd and $N\geq 2$, then $\Phi^{(3)}$ restricts to a bijection
    \begin{equation}
        \Phi^{(3)}: {}^\dagger C_1(2P)\times {}^\dagger C_{N-1}(P)\xrightarrow{\sim} {}^\dagger C_N(P).
    \end{equation}
\end{lemma}

\begin{proof}
    Let us first show that the map is well-defined.
    Given $(n_1,n_2,n_3)\in {}^\dagger C_1(2P)$ and $(b_1,b_2,b_3)\in {}^\dagger C_{N-1}(P)$, we need to show that $(\Phi(n_i,b_i))_{i=1,2,3}\in {}^\dagger C_N(P)$.

    Let us check that \eqref{eq: lower} is satisfied for $M<N$.
    In this case, combining \eqref{eq: am} and \eqref{eq: Phinb}, we see that $\Phi(n_i,b_i)_M$ is either $[b_i]_M$ or $P^M-1-[b_i]_M$.
    Therefore, \eqref{eq: lower} is satisfied for $(\Phi(n_i,b_i))_{i=1,2,3}$ since it is satisfied for $(b_1,b_2,b_3)$. 
    

    We now check that \eqref{eq: upper} is also satisfied for $(\Phi(n_i,b_i))_{i=1,2,3}$.
Note that  $2\Phi(n_i,b_i)+1=P^{N-1}n_i+w_i$, where $w_i=2b_i+1$ if $n_i$ is even, and $w_i=P^{N-1}-(2b_i+1)$ if $n_i$ is odd.
By Lemma \ref{eqn: four ineq}, the condition \eqref{eq: upper} for $(\Phi(n_i,b_i))_{i=1,2,3}$ is satisfied if and only if we have 
\begin{align}\label{eq: P^N-1 four ineq}
\begin{split}
    P^{N-1}(2P-2-\sum_{i=1}^3 n_i)+\Big(2P^{N-1}-\sum_{i=1}^3 w_i\Big)&>0,\\
    P^{N-1}(-n_i+\sum_{j\neq i}n_j)+\Big(-w_i+\sum_{j\neq i}w_j\Big)&>0,\quad i=1,2,3.
\end{split}
\end{align}


To continue the argument, we set up the following notation.
Let $H_0=2P^{N-1}-\sum_{i=1}^3(2b_i+1)$, and let $H_i=-(2b_i+1)+\sum_{j\neq i}(2b_j+1)$. 
By Lemma \ref{eqn: four ineq}, the condition that $(b_1,b_2,b_3)\in {}^\dagger C_{N-1}(P)$ is equivalent to the positivity of the $H_i$'s.
We first observe that $H_i<2P^{N-1}$ for $i=0,1,2,3$.  Indeed, $H_0=2P^{N-1}-\sum_{i=1}^3(2b_i+1)<2P^{N-1}$.
For $i=1,2,3$, we have that

\[H_i=2(-b_i+\sum_{j\neq i}b_j)+1\leq 1+2\sum_{j\neq i}b_j\leq 1+2(P^{N-1}-2)<2P^{N-1}.\]

Now consider the four inequalities in \eqref{eq: P^N-1 four ineq}.
All the terms in the first brackets are non-negative, since
$(n_1,n_2,n_3)\in{}^\dagger C_1(2P)$.
Moreover, all four first-bracket terms
have the same parity as $n_1+n_2+n_3$.

The following table records the four second-bracket terms in
\eqref{eq: P^N-1 four ineq}.  In the rows involving $i,j,k$, the indices
$i,j,k$ are distinct elements of $\{1,2,3\}$, and the four entries are listed
in the respective order.
\[
\begin{array}{c|c}
\text{parity pattern} & \text{four second-bracket terms} \\ \hline
n_1,n_2,n_3 \text{ all even}
& (H_0,H_1,H_2,H_3)\\
n_i,n_j \text{ odd and } n_k \text{ even}
& (H_k,H_j,H_i,H_0)\\
n_i \text{ odd and } n_j,n_k \text{ even}
& (P^{N-1}-H_i,\ P^{N-1}-H_0,\ P^{N-1}-H_k,\ P^{N-1}-H_j)\\
n_1,n_2,n_3 \text{ all odd}
& (P^{N-1}-H_0,\ P^{N-1}-H_1,\ P^{N-1}-H_2,\ P^{N-1}-H_3).
\end{array}
\]
Thus, when $n_1+n_2+n_3$ is even, the second-bracket terms are a permutation
of $(H_i)_{i=0,1,2,3}$; when $n_1+n_2+n_3$ is odd, they are a permutation of
$(P^{N-1}-H_i)_{i=0,1,2,3}$.

If $n_1+n_2+n_3$ is even, then every second-bracket term is one of the positive
numbers $H_i$, so all four inequalities in \eqref{eq: P^N-1 four ineq} hold.
If $n_1+n_2+n_3$ is odd, then every first-bracket term is at least $1$, while
every second-bracket term is of the form $P^{N-1}-H_i>-P^{N-1}$, since
$H_i<2P^{N-1}$.  Hence the corresponding left hand side is strictly positive
in this case as well.

Above shows that $\Phi^{(3)}$ sends ${}^\dagger C_1(2P)\times {}^\dagger C_{N-1}(P)$ to ${}^\dagger C_{N}(P)$. We now construct its inverse.


Take an arbitrary element $(a_1,a_2,a_3)\in{}^\dagger C_N(P)$.
  By \eqref{eq: am} and \eqref{eq: lower} for \(M=N-1\), we can
choose
$b_i\in\{[a_i]_{N-1},\,P^{N-1}-1-[a_i]_{N-1}\}$
such that $(b_1,b_2,b_3)$ satisfies \eqref{eq: lower} with $M=N-1$.

We claim that $(b_1,b_2,b_3)\in{}^\dagger C_{N-1}(P)$.  Indeed, it remains
only to check \eqref{eq: am} and \eqref{eq: lower} for $M<N-1$. 
For such $M$,  we have $\{[b_i]_M,\,P^M-1-[b_i]_M\}
=
\{[a_i]_M,\,P^M-1-[a_i]_M\}$.
Therefore the choices that work for $(a_1,a_2,a_3)$ also work for
$(b_1,b_2,b_3)$.

Moreover, the choice of each $b_i$ is unique.  To see this, note that
\eqref{eq: upper} implies that every coordinate of an element of ${}^\dagger C_{N-1}(P)$ is at most $\frac{P^{N-1}-3}{2}$.
Indeed, if $(x_1,x_2,x_3)\in{}^\dagger C_{N-1}(P)$, then
$x_i\leq \sum_{j\neq i}x_j$ and $\sum_i x_i\leq P^{N-1}-2$, hence $2x_i\leq P^{N-1}-2$.  Since $P^{N-1}$ is odd, this gives the desired
bound.  
Therefore, among the two numbers $[a_i]_{N-1}$ and $P^{N-1}-1-[a_i]_{N-1}$,
at most one can be a coordinate of an element of
${}^\dagger C_{N-1}(P)$.

We can now define $n_i$ so that $a_i=\Phi(n_i,b_i)$.
Namely, let
\[
        n_i=
        \begin{cases}
        2\frac{a_i-[a_i]_{N-1}}{P^{N-1}},
        & b_i=[a_i]_{N-1},\\[6pt]
        2\frac{a_i-[a_i]_{N-1}}{P^{N-1}}+1,
        & b_i=P^{N-1}-1-[a_i]_{N-1}.
        \end{cases}
\]

It remains to show that $(n_1,n_2,n_3)\in{}^\dagger C_1(2P)$, i.e., the first brackets in the four inequalities in \eqref{eq: P^N-1 four ineq} are all non-negative.  
The argument is the reverse of the argument for well-definedness.

Since
$(a_1,a_2,a_3)\in{}^\dagger C_N(P)$, Lemma \ref{eqn: four ineq} gives the positivity of the four inequalities in \eqref{eq: P^N-1 four ineq}.  Suppose,
for contradiction, that one of the terms in the first brackets in
\eqref{eq: P^N-1 four ineq} is negative.

If $n_1+n_2+n_3$ is even, then every term in the first brackets is even, so
the negative first-bracket terms are $\leq -2$.  We have already shown that $H_i<2P^{N-1}$.
Thus the corresponding left hand side in \eqref{eq: P^N-1 four ineq} is
strictly negative, a contradiction.

If $n_1+n_2+n_3$ is odd, then every term in the first brackets is odd, so the
negative first-bracket terms are $\leq -1$.  In this case, the corresponding second bracket
is  $P^{N-1}-H_i<P^{N-1}$. Again the corresponding left hand side in
\eqref{eq: P^N-1 four ineq} is strictly negative, a contradiction.

Therefore, we have found the unique $(n_1,n_2,n_3)\in {}^\dagger C_1(2P)$ and $(b_1,b_2,b_3)\in {}^\dagger C_{N-1}(P)$, which defines $\Phi^{(3),-1}(a_1,a_2,a_3)$.
\end{proof}

\end{document}
